\kbvcontinued ದಿಂದ “ಅನ್ನು” ಎಂಬ ಪ್ರತ್ಯಯ ಲೋಪವಾಗಿ, ಆಟ + ಗಾರ = ಆಟಗಾರ ಎಂದಾಗುವುದು. ಇಲ್ಲಿ ಬಳಸಿರುವ “ಗಾರ” ಪ್ರತ್ಯಯವನ್ನು ತದ್ಧಿತಪ್ರತ್ಯಯ ವೆನ್ನುತ್ತೇವೆ. ಈ ರೀತಿಯಾದ ಪ್ರತ್ಯಯಗಳನ್ನು ಅಂತ್ಯದಲ್ಲಿ ಉಳ್ಳ ಶಬ್ದರೂಪಗಳೇ ತದ್ಧಿತಾಂತ (ತದ್ಧಿತ + ಅಂತ)ಗಳು.\par

\kbvsechead{ಈ ಉದಾಹರಣೆಗಳನ್ನು ಗಮನಿಸಿ:}

\vspace{2.5pt plus 3pt minus 0.5pt}\noindent\hspace*{1.2em}\kbvnum{(1)}\hspace{0.55em}ಲಂಚ + (ಅನ್ನು) + ಕೋರ = ಲಂಚಕೋರ {[}ಲಂಚವನ್ನು ಕೋರುವವನು ಲಂಚಕೋರ{]}\par

\vspace{2.5pt plus 3pt minus 0.5pt}\noindent\hspace*{1.2em}\kbvnum{(2)}\hspace{0.55em}ಚಿತ್ರ + (ಅನ್ನು) + ಕಾರ = ಚಿತ್ರಕಾರ {[}ಚಿತ್ರವನ್ನು ಬರೆಯುವವನು ಚಿತ್ರಕಾರ{]}\par

\vspace{2.5pt plus 3pt minus 0.5pt}\noindent\hspace*{1.2em}\kbvnum{(3)}\hspace{0.55em}ಗಾಣ + (ಅನ್ನು) + ಇಗ = ಗಾಣಿಗ {[}ಗಾಣವನ್ನು ಆಡಿಸುವವನು ಗಾಣಿಗ{]}\par

\vspace{2.5pt plus 3pt minus 0.5pt}\noindent\hspace*{1.2em}\kbvnum{(4)}\hspace{0.55em}ಮಡಿ + (ಅನ್ನು) + ವಂತ = ಮಡಿವಂತ {[}ಮಡಿಯನ್ನು ಆಚರಿಸುವವನು ಮಡಿವಂತ{]}\par

\vspace{2.5pt plus 3pt minus 0.5pt}\noindent\hspace*{1.2em}\kbvnum{(5)}\hspace{0.55em}ಕೌರ + (ಅನ್ನು) + ಇಕ = ಕೌರಿಕ {[}ಕೌರವನ್ನು ಮಾಡುವವನು ಕೌರಿಕ{]}\par

\indent ಮೇಲೆ ತಿಳಿಸಿದ ಉದಾಹರಣೆಗಳಲ್ಲಿ ಕೋರ, ಕಾರ, ಇಗ, ವಂತ ಇಳಿತ್ತು. ಈ ಪ್ರತ್ಯಯಗಳು,\par

\begin{kbvrule}
\noindent\kbvbullet\hspace{0.5em}ತದ್ಧಿತ ಪ್ರತ್ಯಯಗಳು ನಾಮಪದಗಳ ಮೇಲೆ ಬಂದು ಹೊಸ ಶಬ್ದಗಳಾಗುವುದು. ಪ್ರತ್ಯಯಗಳಿಗೆ ತದ್ಧಿತಾಂತಗಳೆಂದು ಹೆಸರು. ತದ್ಧಿತಾಂತಗಳನ್ನು (1) ತದ್ಧಿತಾಂತ ನಾಮ, (2) ತದ್ಧಿತಾಂತ ಭಾವನಾಮ, (3) ತದ್ಧಿತಾಂತಾವ್ಯಯ ಎಂದು ಮೂರು ಬಗೆಯಾಗಿ ವಿಂಗಡಿಸಬಹುದು.
\end{kbvrule}

\vspace{4pt plus 4pt minus 1pt}\noindent{\kbvsecfont\fontsize{11pt}{14.5pt}\selectfont\color{kbvaccentdark} (1)\hspace{0.45em}ತದ್ಧಿತಾಂತನಾಮ:}\hspace{0.5em}ನಾಮಪದಗಳ ಮೇಲೆ ವಿವಿಧ ಅರ್ಥಗಳಲ್ಲಿ ಗಾರ, ಕಾರ, ಇಗ, ವಂತ ಮೊದಲಾದ ತದ್ಧಿತ ಪ್ರತ್ಯಯಗಳು ಸೇರಿ ಆಗುವ ಹೊಸ ಶಬ್ದಗಳಿಗೆ ತದ್ಧಿತಾಂತ ನಾಮಗಳೆಂದು ಹೆಸರು.\par
